For every fixed $c>0$ and all sufficiently large $n$, an $n$-uniform family $\mathcal F$ with $|\mathcal F|\le c2^n$ has at least $\eta_c2^{|X|}$ proper two-colourings of $X=\bigcup\mathcal F$, for some $\eta_c>0$. Taking $B$ to be one colour class gives exactly the sets requested: each edge meets $B$ and its complement.

We reconstruct Koishi Chan's adaptive-exposure proof at \url{https://www.erdosproblems.com/forum/thread/1027}, keeping the factors in the residual weights explicit. The external input is Beck's nonuniform property-B theorem: for every $C>0$ there is an integer $R$ such that every finite hypergraph with minimum edge size at least $R$ and
\[
 \sum_e2^{-|e|}\le C
\]
has a proper two-colouring. This is the consequence $f(R)\to\infty$ proved by Beck, and strengthened in Linyuan Lu's Theorem 1, \emph{On a problem of Erd\H{o}s and Lov\'asz on coloring non-uniform hypergraphs}, \url{https://people.math.sc.edu/lu/papers/propertyB.pdf}. The nonuniform theorem is assumed; the positive proportion of colourings is proved below.
